\chapter{The Marxist Perspective: A. R. Desai}

\section{A. R. Desai: Life \& Influences}

A. R. Desai is the most important Marxist sociologist of India; his work brought a \textbf{radical turn} to Indian sociology. The Bombay school had two legacies --- started by Ghurye, finished by Desai (D. N. Majumdar's point). \textbf{Sujata Patel} is to Desai what Carol Upadhyay is to Ghurye: she says Desai is ``first and foremost a Marxist, then a sociologist, then a human being''.

\begin{itemize}
\item Influences: (1) the \textbf{nationalist movement}; (2) \textbf{Marxism}; (3) his father \textbf{Ramlal Desai} (a civil servant, storyteller, novelist and Gandhian).
\item In \textbf{1934} he joined the Communist Party of India and was \textbf{rusticated} from Baroda University for radical activism; he then studied under Ghurye in Bombay.
\end{itemize}

\begin{KeyConcept}
Marx: ``Philosophers have only interpreted the world; the point is to change it.'' Desai subscribed to this --- he wanted to change society, not just understand it.
\end{KeyConcept}

His famous book \emph{Social Background of Indian Nationalism} (1948); an earlier book was \emph{Gandhi X-Ray} (critical of Gandhi --- his father loved Gandhi, he did not).

\begin{QuickRevision}
\begin{itemize}
\item A. R. Desai = most important Marxist sociologist; radical turn in Indian sociology.
\item Influences: nationalist movement, Marxism, father Ramlal Desai.
\item 1934 joined CPI, rusticated from Baroda; studied under Ghurye.
\item Book: Social Background of Indian Nationalism (1948); Gandhi X-Ray.
\end{itemize}
\end{QuickRevision}

\section{Feudalism, Colonialism \& Capitalism}

Desai reads India through the \textbf{mode of production}: pre-British India was \textbf{feudal}; colonialism introduced \textbf{capitalism}; after independence India was supposed to follow socialism but never did --- the Indian state became the \textbf{guardian of capitalism}.

\begin{KeyConcept}
Features of Indian feudalism: (1) \textbf{self-sufficient village economy}; (2) \textbf{caste structure} --- ``caste is the steel frame of Hinduism'' (like bureaucracy is the steel frame of India, per Sardar Patel), with guilds (shrenis) based on caste/occupation; (3) \textbf{communal ownership of land} (no private property); (4) \textbf{revenue collector as representative of the monarch}.
\end{KeyConcept}

Desai read the economist \textbf{R. P. Dutt} (who wrote \emph{India Today}) and \textbf{Dadabhai Naoroji} (drain of wealth).

\begin{QuickRevision}
\begin{itemize}
\item Feudal (pre-British) to capitalist (colonialism) to state as guardian of capitalism (post-independence).
\item Feudalism: self-sufficient village, caste = steel frame of Hinduism, communal land, revenue collector.
\end{itemize}
\end{QuickRevision}

\section{The Critique of the Indian State}

\begin{itemize}
\item \emph{State and Society} (1976) contains the essay \textbf{``The Myth of Welfare State''} --- the Indian state claims to be a welfare state but serves the capitalist class, not the proletariat.
\item \emph{India's Path of Development} (1978/1980) reviewed 30 years (1948--77) and found increasing class inequality and poverty.
\item Indian National Congress was an ally of the capitalist class, a ``Brahmin-baniya party'' (Periyar's view; the self-respect movement; anti-Brahminical Dravidian movement leading to DMK/AIADMK).
\end{itemize}

\begin{QuickRevision}
\begin{itemize}
\item Myth of Welfare State (1976): state serves capitalist class.
\item India's Path of Development (1978): increasing inequality.
\item Congress = ally of capitalists, ``Brahmin-baniya party'' (Periyar).
\end{itemize}
\end{QuickRevision}

\section{State, Welfarism \& Planning}

Marxism helps understand \textbf{state formation, its ideology and the assessment of welfarism}. The first tool of the Indian state was \textbf{planning} --- but plans were financed by \textbf{taxing the common people}, not the capitalist class. This reduced purchasing power and consumerism, so development happened for the capitalist, never the nation.

\begin{itemize}
\item Industries were controlled by old caste-privileged families (Birla, Tata --- from British times); old banias became new financial/commercial capitalists.
\item \textbf{Land reforms} (intended for the landless) actually created a class more dependent on the state, helping only landholders.
\end{itemize}

\begin{CriticalPoint}
Instruments that seemed socialist (planning, cooperatives, land reform, regional rural banks, mixed economy) were, for Desai, lip service --- serving the capitalist class and the nexus of rural-urban bourgeoisie. The Green Revolution's unanticipated consequence was the creation of \textbf{rich peasants}.
\end{CriticalPoint}

\begin{QuickRevision}
\begin{itemize}
\item Planning taxed the masses, reducing purchasing power; served the capitalist class.
\item Land reform helped landholders; Green Revolution created rich peasants.
\end{itemize}
\end{QuickRevision}

\section{Society as Emerging, Developing \& Declining}

Marxism explains society as \textbf{emerging, developing, subsequently declining, and ultimately either emerging into a qualitatively new society or disintegrating} --- through the inherent contradictions of each mode of production (dialectics is the motor of history).

\begin{QuickRevision}
\begin{itemize}
\item Marxism: society emerges, develops, declines, transforms through inherent contradictions (dialectics = motor of history).
\end{itemize}
\end{QuickRevision}

\section{Nationalism as an Inherent Contradiction of Colonialism}

Nationalism was the \textbf{inherent contradiction of colonialism}: colonialism created the conditions (modern education, modern means of communication) that produced nationalist consciousness, which then turned against colonialism. Desai: ``colonialism gave way to nationalism, and this nationalism, although made by colonialism, turned against colonialism itself''.

\begin{itemize}
\item The \textbf{moderates} lived under \textbf{false consciousness} (they believed the British were modernising India); the \textbf{extremists} exposed the truth (development of underdevelopment).
\end{itemize}

\begin{CriticalPoint}
Desai on Gandhi: Gandhi borrowed mass-line politics from \textbf{Lenin} (mass-based movement, Bolshevik revolution). But Gandhi's approach was actually \textbf{class collaborationist}, not mass-based --- he was a ``friend of industrialists''. In the Ahmedabad mill strike (1918) Gandhi only fought for a wage increase and preached \textbf{trusteeship}, never touching structural exploitation. Trusteeship (donating the surplus) merely \textbf{opiated the masses}.
\end{CriticalPoint}

\begin{QuickRevision}
\begin{itemize}
\item Nationalism = inherent contradiction of colonialism.
\item Moderates = false consciousness; extremists exposed the truth.
\item Gandhi: class collaborationist, trusteeship opiated the masses (Ahmedabad mill strike 1918).
\end{itemize}
\end{QuickRevision}

\section{The Significance of the Marxist Perspective}

\begin{itemize}
\item Marxist perspective was a \textbf{challenge to the status quoist} structural-functional and Indological approaches.
\item It is the only approach that can explain \textbf{social change} (structural functionalism can only explain maintenance).
\item It draws on \textbf{history, political economy and anthropology} --- Desai opposed rigid disciplinary boundaries.
\item It understands India through the \textbf{economic base}, whereas Indology and structural functionalism understand India through culture, religion, family, caste.
\end{itemize}

\begin{KeyConcept}
Micro (field study, Srinivas) vs \textbf{macro} (book view, Ghurye) vs \textbf{mezzo} (middle path, Desai).
\end{KeyConcept}

\begin{CriticalPoint}
Critique of Desai: (1) not substantiated empirically --- clubbing 2000 years as ``feudal'', claiming land was communally owned (Gupta land grants contradict this); (2) not objective --- projecting the Indian state purely as a capitalist instrument, ignoring welfare.
\end{CriticalPoint}

\begin{QuickRevision}
\begin{itemize}
\item Marxist perspective = challenge to status quo; explains social change; draws on history/political economy/anthropology.
\item Macro (Ghurye), micro (Srinivas), mezzo (Desai).
\item Critique: not empirical, not objective.
\end{itemize}
\end{QuickRevision}
